\section{The Optimal Series \texorpdfstring{$18\cdot 2^t+1$}{18·2ᵗ+1}}\label{sec:series}

Since the base configuration of Theorem~\ref{thm:main-seed} already satisfies
the hypotheses of the affine iterative step
(Proposition~\ref{prop:iterative-step}), the construction of
\cite{bartholdi2007} applies directly and yields the following.

\begin{theorem}[Straight-line optimal series for $18\cdot 2^t+1$ lines]\label{thm:straight-series}
For every integer $t\ge 0$, there exists a simple affine arrangement of
$n=18\cdot 2^t+1$ straight lines attaining the affine upper
bound~\eqref{eq:upper-bound}:
\[
  a_3
  =\left\lfloor \frac{n(n-2)}{3}\right\rfloor
  =108\cdot 4^t-1.
\]
\end{theorem}

\begin{proof}
By Theorem~\ref{thm:main-seed}, the base configuration $\cA_0$ is a simple
affine arrangement of $19$ straight lines with $a_3(\cA_0)=107$ bounded
triangles and $Y_0$ touching exactly $17$ of them. By
Lemma~\ref{lem:bbl-match}, the realization holds for all
$\eps\in(0,\SeedNineteenEpsUpper)$; in particular, $\eps$ can be chosen
arbitrarily small, so the iteration condition of
Remark~\ref{rem:iteration-condition} is satisfied and
Proposition~\ref{prop:iterative-step} can be applied repeatedly. Writing
$n_t=18\cdot 2^t$, each step produces a simple affine arrangement $\cA_{t+1}$
of $2n_t+1$ straight lines with $n_t^2$ more bounded triangles than $\cA_t$.
Therefore
\[
  a_3(\cA_t)
  =107+\sum_{s=0}^{t-1}(18\cdot 2^s)^2
  =107+324\cdot\frac{4^t-1}{3}
  =108\cdot 4^t-1
  =\frac{(n_t+1)(n_t-1)-2}{3},
\]
which is the upper bound~\eqref{eq:upper-bound} for $n=n_t+1$.
\end{proof}
